\section*{Capítulo \ref{ch:b3:dna-repair} --- Estabilidad del genoma: daño, reparación y recombinación del ADN}

\begin{solution}{exo:b3:dna-repair:1}
Uracilo: espontánea (\omterm{def:b3:dna-repair:lesions}{desaminación} de la C), \omterm{def:b3:dna-repair:excision}{reparación por escisión de
bases} por la uracilo-ADN-glicosilasa. Dímero de timina: ambiental (UV),
\omterm{def:b3:dna-repair:excision}{reparación por escisión de nucleótidos}. \omterm{def:b3:dna-repair:lesions}{8-oxoguanina}: espontánea
(oxidación), \omterm{def:b3:dna-repair:excision}{reparación por escisión de bases} por la \omterm{def:b3:dna-repair:excision}{glicosilasa} de la
8-oxoG. Mal apareamiento G--T detrás de la horquilla: error de
replicación, \omterm{prop:b3:dna-repair:fidelity}{reparación de errores de apareamiento}. \omterm{def:b3:dna-repair:lesions}{Rotura de doble hebra}
en G1: ambiental (radiación) o espontánea; \omterm{def:b3:dna-repair:dsb}{unión de extremos no
homólogos}, ya que no hay cromátida hermana disponible.
\end{solution}

\begin{solution}{exo:b3:dna-repair:2}
La escisión de nucleótidos reconoce una propiedad física común a todas
las lesiones voluminosas --- la deformación de la hélice, o una
ARN-polimerasa detenida --- y retira un parche que contiene lo que la
causara. Las lesiones pequeñas (uracilo, 8-oxoG, bases alquiladas) apenas
deforman la hélice, de modo que han de reconocerse por su química, una
\omterm{def:b3:dna-repair:excision}{glicosilasa} por cada tipo de base errónea.
\end{solution}

\begin{solution}{exo:b3:dna-repair:3}
Un mal apareamiento consta de dos bases normales, y solo la de la hebra
nueva es la equivocada; el sistema tiene que saber cuál es la hebra
nueva. \emph{E.~coli} marca la hebra parental con grupos metilo en los
GATC, que a la hebra nueva le faltan durante unos minutos; MutH muerde la
hebra no metilada. En un mutante \emph{dam} ninguna hebra está metilada:
MutH muerde cualquiera de las dos, la reparación quita la base correcta
tan a menudo como la equivocada --- y fija así como mutaciones la mitad
de los errores (una mutadora) --- y las mordeduras en ambas hebras en
sitios GATC cercanos producen \omterm{def:b3:dna-repair:lesions}{roturas de doble hebra}.
\end{solution}

\begin{solution}{exo:b3:dna-repair:4}
Unas $35\times 2 = 70$ \omterm{def:b3:dna-repair:lesions}{roturas de doble hebra}. Al cabo de una hora la
mayoría están unidas: con una semivida cercana a \qty{30}{min}, unos
$70/4 \approx 18$ focos. A las seis horas casi todas se han reparado
($70/2^{12} < 1$), y solo quedan unas pocas roturas persistentes y
complejas.
\end{solution}

\begin{solution}{exo:b3:dna-repair:5}
(a) $10^{-5}\times 10^{-2}\times 10^{-3} = 10^{-10}$; $6.4\times
10^{9}\times 10^{-10} = 0.64$ mutaciones por división. (b) Sin
\omterm{prop:b3:dna-repair:fidelity}{reparación de errores de apareamiento}, $10^{-7}$; $640$ por división.
(c) Sin \omterm{prop:b3:dna-repair:fidelity}{corrección de pruebas}, $10^{-8}$; $64$ por división. La mutante
en \omterm{prop:b3:dna-repair:fidelity}{reparación de errores de apareamiento} es la más peligrosa, mil veces
por encima de lo normal.
\end{solution}

\begin{solution}{exo:b3:dna-repair:6}
$k = \ln 2/\qty{0.25}{h} = \qty{2.8}{h^{-1}}$; $D = 10^{4}/24 =
\qty{420}{h^{-1}}$; $L^{*} = 420/2.8 \approx 150$ sitios abásicos
presentes en cada instante. Con la reparación diez veces más lenta,
$L^{*} \approx \num{1500}$. La horquilla pasa una vez por cada locus y
encuentra allí las lesiones presentes en ese momento; promediado sobre el
genoma encuentra unos $L^{*}$ de ellos --- unos $150$ en condiciones
normales y $\num{1500}$ con el fármaco.
\end{solution}

\begin{solution}{exo:b3:dna-repair:7}
$k_{\text{NHEJ}} = \ln 2/0.5 = \qty{1.39}{h^{-1}}$, $k_{\text{HR}} =
\ln 2/4 = \qty{0.17}{h^{-1}}$; fracción de HR, $0.17/1.56 = 0.11$;
semivida, $\ln 2/1.56 = \qty{27}{min}$. Una horquilla colapsada da una
rotura de \emph{un solo extremo}: no hay un segundo extremo que la NHEJ
pueda unir, de modo que solo la recombinación con la hermana puede
restaurar la horquilla --- y es, en todo caso, la única vía exacta.
\end{solution}

\begin{solution}{exo:b3:dna-repair:8}
(a) Lynch: microsatélites inestables (las longitudes difieren entre
células del tumor), cánceres colorrectal, de endometrio, gástrico y de
ovario, y respuesta normal al sol. (b) XP-A: microsatélites estables,
quemaduras extremas con una exposición mínima, cánceres de piel y de ojo
en la infancia y, en este grupo, degeneración neurológica progresiva.
(c) XP-V: microsatélites estables, reparación por escisión normal, de
modo que la quemadura es más leve, pero los dímeros que escapan a la
escisión los sortean polimerasas propensas al error --- cánceres de piel,
de aparición más tardía que en el grupo~A.
\end{solution}

\begin{solution}{exo:b3:dna-repair:9}
Células hepáticas: la Cre ha escindido el exón~3 de los dos alelos (como
un círculo que se pierde), de modo que el gen está inactivo en todos los
hepatocitos desde el nacimiento --- una supresión específica del hígado.
Neuronas: sin Cre, los alelos flanqueados por lox están intactos y son
funcionales. Orientaciones opuestas: la Cre invierte el exón y lo vuelve
a invertir indefinidamente, de modo que en cualquier momento cerca de la
mitad de las células llevan el exón invertido, no funcional --- un
mosaico, no una deleción limpia.
\end{solution}

\begin{solution}{exo:b3:dna-repair:10}
A medida que LexA se corta, su concentración baja poco a poco. Los
operadores que unen LexA con poca fuerza quedan libres con una caída
pequeña --- \emph{uvrA}, \emph{uvrB}, \emph{recA} ---, de modo que la
reparación exacta empieza primero; los operadores de \emph{sulA} y de las
polimerasas propensas al error unen LexA con fuerza y solo quedan libres
cuando ya casi no queda, es decir, cuando el daño ha persistido mucho
tiempo. Una bacteria sin las polimerasas propensas al error moriría en
las horquillas bloqueadas por lesiones que no pudiera sortear y, sin
mutagénesis inducida, desarrollaría resistencia al antibiótico mucho más
despacio --- una desventaja para la bacteria y una ventaja para el
paciente; por eso se estudian inhibidores de la \omterm{prop:b3:dna-repair:sos}{respuesta SOS}.
\end{solution}

\begin{solution}{exo:b3:dna-repair:11}
Cadena: PARP inhibida $\to$ las roturas de hebra sencilla persisten
$\to$ una horquilla de replicación convierte cada una en una \omterm{def:b3:dna-repair:lesions}{rotura de
doble hebra} de un solo extremo $\to$ su reparación exige \omterm{def:b3:dna-repair:dsb}{recombinación
homóloga} $\to$ la célula tumoral, sin \omterm{def:b3:dna-repair:dsb}{BRCA}, no puede $\to$ uniones
erróneas, pérdida de cromosomas, muerte. Las células normales, con un
alelo funcional, recombinan y sobreviven. Resistencia: (1) una segunda
mutación de \emph{\omterm{def:b3:dna-repair:dsb}{BRCA}} que restaure el marco de lectura y dé una
proteína que funcione --- se detecta secuenciando el tumor o el ADN
tumoral circulante en sangre, y por el regreso de los focos de \omterm{def:b3:dna-repair:dsb}{Rad51} tras
la irradiación; (2) la pérdida de las proteínas que bloquean la resección
de extremos (53BP1--escudina), que restaura en parte la recombinación sin
BRCA1 --- se detecta por su ausencia en la tinción y por los focos de
\omterm{def:b3:dna-repair:dsb}{Rad51} recuperados; también mutaciones de PARP1 que impiden que quede
atrapada en el ADN, o bombas de expulsión del fármaco, detectadas por
secuenciación y por expresión respectivamente.
\end{solution}

\begin{solution}{exo:b3:dna-repair:12}
Dos dúplex homólogos, uno con $A$ y el otro con $a$, intercambian hebras
sencillas a lo largo del marcador: cada uno tiene ahora una hebra de cada
progenitor en ese tramo --- un heterodúplex con un mal apareamiento
$A/a$. Si la \omterm{prop:b3:dna-repair:fidelity}{reparación de errores de apareamiento} convierte el mal
apareamiento del dúplex invadido en $A$, tres cromátidas llevan $A$ y una
$a$: $3{:}1$; convertido en $a$: $1{:}3$; si se deja sin reparar, la
cromátida segrega $A$ y $a$ en la primera mitosis posterior a la meiosis
y da dos esporas hijas distintas (segregación posmeiótica, $5{:}3$ en un
asca de ocho esporas). El tramo de heterodúplex es solo de unos
centenares a unos miles de pares de bases alrededor de la rotura
iniciadora, de modo que solo pueden convertirse los marcadores que caen
dentro y estos son, por construcción, vecinos del entrecruzamiento.
\end{solution}

\begin{solution}{pb:b3:dna-repair:1}
\textbf{1.} $10^{4} + 400 + 10^{4} \approx 2\times 10^{4}$ lesiones al
día, $7\times 10^{6}$ al año --- aproximadamente una por kilobase de
genoma y año.
\textbf{2.} $10^{-5}\times 10^{-2}\times 10^{-3} = 10^{-10}$ por bp;
$6.4\times 10^{9}\times 10^{-10} = 0.64$ mutaciones por división.
\textbf{3.} $0.64\times 0.015 \approx 0.01$ mutaciones codificantes por
división; a lo largo de $50$ divisiones, unas $0.5$.
\textbf{4.} $10^{-7}$ por bp, $640$ mutaciones por división. Células
con resbalón: sin reparación,
$10^{-4}\times 40\times 10^{9} = 4\times 10^{6}$; con ella,
$10^{-7}\times 40\times 10^{9} = 4000$ --- una diferencia de mil veces,
visible como \omterm{def:b3:dna-repair:mmr}{inestabilidad de microsatélites}.
\textbf{5.} El uracilo es ajeno al ADN, lo reconoce la
uracilo-ADN-glicosilasa y se escinde: se repara. La timina procedente de
la 5-metilcitosina es una base normal frente a una G; nada la señala como
errónea fuera de la replicación, y se convierte en una mutación
C$\to$T --- el mecanismo que ha empobrecido el genoma en \omterm{def:b3:chromatin-epigenetics:methylation}{CpG}.
\textbf{6.} La \omterm{def:b3:dna-repair:excision}{glicosilasa} de la 8-oxoG retira la base oxidada mientras
sigue frente a una C (antes de la replicación); la segunda \omterm{def:b3:dna-repair:excision}{glicosilasa}
retira una A mal incorporada frente a la 8-oxoG (después de la
replicación, lo que devuelve una oportunidad de reparar la G); y la
hidrolasa destruye el dGTP oxidado para que no se incorpore frente a una
A. Cuando fallan las tres, la 8-oxoG se aparea con A y la ronda siguiente
fija una transversión G$\cdot$C$\to$T$\cdot$A.
\textbf{7.} Normal, $k = \ln 2/2 = \qty{0.35}{h^{-1}}$; XP, $k = \ln 2/70
= \qty{0.0099}{h^{-1}}$.
\textbf{8.} Normal, $10^{5} e^{-2.77} \approx 6200$; XP, $10^{5}
e^{-0.079} \approx \num{92000}$.
\textbf{9.} Mutaciones, $6.2$ y $92$ por célula; codificantes, $0.09$ y
$1.4$.
\textbf{10.} Con $500$ exposiciones: $47$ mutaciones codificantes
(normal) y $690$ (XP), frente a $0.5$ por la replicación --- el sol
domina la carga de mutaciones de la piel incluso en una persona normal.
\textbf{11.} Secuencia codificante, $9.6\times 10^{7}$ bp; los genes
conductores son $4\times 10^{4}/9.6\times 10^{7} = 4.2\times 10^{-4}$ de
ella. Impactos medios, $690\times 4.2\times 10^{-4} = 0.29$;
$P(\ge 1) = 1 - e^{-0.29} =
0.25$. De $10^{9}$ células, $2.5\times 10^{8}$ llevan una mutación
conductora (niño normal: media $0.02$, \qty{2}{\%}).
\textbf{12.} Un dímero no es una mutación; solo se convierte en una
cuando una horquilla lo copia mediante un salto propenso al error, de
modo que las que mutan son las células que se dividen con muchos dímeros
sin reparar. La luz ultravioleta solo llega a la piel y al ojo; los
tejidos internos no reciben dímeros y, en la XP, son casi normales.
\textbf{13.} $70$ roturas; al cabo de \qty{1}{h} solo con NHEJ, $70\times
2^{-2} \approx 18$ focos.
\textbf{14.} $k_{\text{NHEJ}} = \qty{1.39}{h^{-1}}$, $k_{\text{HR}} =
\qty{0.17}{h^{-1}}$; fracción de HR, $0.17/1.56 = 0.11$.
\textbf{15.} Semivida, $\ln 2/1.56 = \qty{27}{min}$; al cabo de
\qty{1}{h}, $70\,e^{-1.56} \approx 15$ roturas.
\textbf{16.} $70\times 0.015\times 0.7 \approx 0.7$ genes estropeados por
célula.
\textbf{17.} $n = 70$: $2415$ pares $\times 3\times 10^{-4} = 0.72$
translocaciones. $n = 7$: $21\times 3\times 10^{-4} = 0.006$. El número
de roturas es $\propto D$, pero el número de \emph{pares} de roturas
simultáneas es $\propto n^{2} \propto D^{2}$.
\textbf{18.} $D = \qty{7}{h^{-1}}$, $L^{*} = 7/1.39 \approx 5$ roturas
presentes a la vez: $10$ pares, $0.003$ translocaciones --- doscientas
veces menos que la misma dosis en un instante. El fraccionamiento deja
que el tejido normal repare entre fracciones y mantiene pocas las
roturas simultáneas.
\textbf{19.} Reparación intacta: $e^{-2/1.5} = 0.26$ a \qty{2}{Gy} y
$e^{-4} = 0.018$ a \qty{6}{Gy}. Sin NHEJ: $e^{-4} = 0.018$ y
$e^{-12} = 6\times 10^{-6}$. $D_{0}$ es la dosis que deja por término
medio una \omterm{def:b3:dna-repair:lesions}{lesión} letal sin reparar por célula (supervivencia $1/e$).
\textbf{20.} El \qty{11}{\%} de las roturas de dos extremos de G2 que la
HR habría reparado con exactitud pasa ahora por la NHEJ, y toda rotura
de horquilla de un solo extremo --- que la NHEJ no puede reparar bien en
absoluto --- se une mal: a lo largo de los años el genoma acumula
deleciones, translocaciones y cambios de número de copias, las cicatrices
características de los tumores \emph{\omterm{def:b3:dna-repair:dsb}{BRCA}}.
\textbf{21.} $10^{4}\times 0.01 = 100$ roturas de un solo extremo al día.
Una rotura de un solo extremo no tiene extremo compañero; la NHEJ solo
puede dejarla o unirla a algún extremo no relacionado, y produce una
translocación o una pérdida.
\textbf{22.} Un alelo funcional de \emph{BRCA2} fabrica proteína
suficiente para la recombinación: las células normales reparan sus
roturas de horquilla y sobreviven al inhibidor.
\textbf{23.} Un segundo desplazamiento del marco aguas abajo del primero
restaura el marco de lectura y da una BRCA2 acortada pero en parte
funcional, de modo que la recombinación vuelve y el fármaco deja de
matar. Se detecta secuenciando \emph{BRCA2} en el tumor resistente o en
el ADN tumoral circulante, y funcionalmente por la reaparición de los
focos de \omterm{def:b3:dna-repair:dsb}{Rad51} tras la irradiación.
\textbf{24.} La muerte necesita que la \omterm{prop:b3:dna-repair:fidelity}{reparación de errores de
apareamiento} actúe sobre las bases alquiladas; un tumor de Lynch, que
carece de ella, es tolerante --- el fármaco no consigue matarlo (esos
tumores se tratan por otros medios, entre ellos la inmunoterapia, a la
que los hace sensibles su elevada carga de mutaciones).
\textbf{25.} Unos \num{92000} dímeros en la horquilla en la célula XP;
$10^{-10}$ mutaciones por par de bases y división; la HR repara el
\qty{11}{\%} de las roturas de G2.
\end{solution}
